\subsubsection*{The abelian theta coefficients}
\label{subsec:toric-twist}
\hypertarget{source-toric-twist}{}

The Selmer determinant controls the signature contribution to the simultaneous twist. The abelian theta values require a separate argument: changing the toric order changes both the coefficients and their period normalization. Let \(E/F\) be the CM extension of the
selected host, with \(F\) totally real, and let \(\varphi\)
be its chosen CM type. Fix a toric label, its integer
\(s\ge1\), the nonnegative harmonic exponents \(\kappa\),
its base conductor \(\mathfrak f\), and its unitary Hecke
character \(\chi\), trivial on the ideles of \(F\).
Its infinity exponents are \(\kappa_i+s\) at the places
of \(\varphi\) in the unitary normalization. At the base
conductor its local conductor is prescribed and nontrivial.
Enlarge the order conductor to
\[
 \mathfrak c=\mathfrak f\prod_{v\in Q}v
\]
in the real
field, at split places prime to the base conductor and to
the coefficient prime \(\ell\).
Let \(\operatorname{Pic}(\mathfrak c)\) denote the
ideal class group of the resulting order.

For a finite idele \(y\) representing a class, let
\(S_y(z)\) be the continuation of the lattice sum on
\(E\cap y\widehat{\mathcal O}_{\mathfrak c}\), with the
harmonic numerator and shifted exponent of the fixed label,
over nonzero vectors modulo \(\mathcal O_F^\times\).
The toric term is
\[
       T_y=\frac{S_y(0)}
          {\chi_f(y)\,|y|_f^{s}},
\]
where the finite modulus \(|y|_f\) is the one in the
unitary Hecke normalization. In the convergent region the
weights are homogeneous under
multiplication of the lattice by an element of \(E^\times\).
The character at infinity and the norm exponent give
the inverse of this change in the finite prefactor.
The completed order units preserve the lattice and
the finite character is trivial on them. The resulting
normalized term therefore depends on the class. These
identities pass to the continued value.

\paragraph{The denominator after changing the order.}
Choose the fixed CM period lattice with basis
\(1,(1+\sqrt{-d})/2\) at the imaginary quadratic subfield,
and choose a twelfth root of its discriminant. Incorporate the
fixed power of this root and the normalization constant
of the harmonic numerator into a nonzero scalar
\(\Omega_*\). In this normalization
\[
 \beta(T)=(T/\Omega_*)^{12}.
\]

The local class-value estimate for lattices with
\(y_w=1\) at every \(w\mid\ell\) has the following form.
Put \(W=\operatorname{wtM}(2s,\kappa)\), the degree of
the period used in the twelfth power. Let \(a_0\) be the
maximum conductor exponent of \(\mathfrak f\) over \(\ell\).
Let \(c_E\) be the largest denominator exponent at
\(\ell\) of the partial zeta values of \(F\) at \(1-2s\),
divided by \(2^{[F:\mathbb Q]}\) before taking the
denominator. Put \(f_0=|I|=[F:\mathbb Q]\) for the size
of the chosen CM type, \(K_\kappa=\sum_{i\in I}\kappa_i\),
and \(W=2sf_0+2K_\kappa\).
Write \(c_\ell=a_0\) if \(a_0>0\) and
\(c_\ell=v_\ell(D_E)/[E:\mathbb Q]\) otherwise. The
ramified allowance is
\[
 n_\ell =
 \begin{cases}
 0,&c_\ell=0,\\
 2sf_0c_\ell+K_\kappa(2c_\ell+1+
       \mathbf 1_{\ell=2}+\tfrac12\mathbf 1_{\ell=d}),&c_\ell>0 .
 \end{cases}
\]
\begin{proposition}[Uniform denominator of the CM class values]\label{prop:toric-class-denominator}
At every embedding into \(\overline{\Q}_\ell\), for a representative
with \(y_w=1\) at all \(w\mid\ell\), the normalized twelfth
power is algebraic and satisfies
\[
    v_\ell\bigl(\beta(S_y(0))\bigr)\ge-12(c_E+n_\ell+Wa_0).
\]
Auxiliary split places disjoint from \(\ell\) may vary.
\end{proposition}
The proof follows in the local jet calculation below. First we pass
from one lattice to the coefficients of the theta element.

\begin{lemma}[Clearing the abelian theta coefficients]\label{lem:toric-common-denominator}
With the CM type, base conductor, critical weight and coefficient
prime fixed, there is an integer \(c\ge0\) such that every coefficient \(a\) of a normalized abelian theta element, at any allowed auxiliary level and after any allowed finite-order ring class twist, satisfies
\[
 v_\ell(a)\ge-c.
\]
\end{lemma}

\begin{proof}
Every ideal class admits a representative that is one
above \(\ell\). Indeed, use weak approximation in \(E\)
to put the components above \(\ell\) in the local
completed order-unit subgroups and then divide by those
units. On such a representative \(\chi_f(y)\) is a unit
at \(\ell\) in the algebraic normalization. This follows
from the formula for its infinity type on principal
ideles, applied after taking a power to clear the ideal
class, and from the roots of unity on the compact unit
group.

The element of \(E^\times\) used in this formula
has valuations zero above \(\ell\), so every conjugate
factor is a unit. The character and norm factors of such a representative have
\[
 v_\ell(\chi_f(y))=0,\qquad v_\ell(|y|_f)=0,
\]
where \(|y|_f\) denotes the finite norm factor in this normalization.
The same bound thus holds for \(\beta(T_y)\) on all
classes. If a finite order ring class character is added,
its values again have valuation zero. If the separately cleared summands \(a_y\) have nonnegative valuation, the only estimate required for their sum is
\[
 v_\ell\!\left(\sum_y a_y\right)\ge\min_y v_\ell(a_y)\ge0.
\]
Thus the number of classes of the varying order does not change the clearing power.

For the abelian theta module the coefficient at an Artin
label is obtained from the terms \(T_y\) in its fiber,
multiplied by the same nonzero normalization scalar
specified by the toric construction. The fiber map
uses the inverse ideal class. Choose the complex
embedding in the algebraic closure of the row fraction
field as in the toric comparison. Taking twelfth roots gives the individual integral bound
\[
 v_\ell(T_y/\Omega_*)\ge -N\qquad(N\in\mathbb Z,\ N\ge c_E+n_\ell+Wa_0).
\]

One more fixed power clears the product of
\(\Omega_*\) and the normalization scalar in that
algebraic closure. To see this for any algebraic element,
take a monic polynomial over the fraction field and
choose the power large enough that the rescaled monic
polynomial has integral coefficients. Each coefficient
of the theta module is already in the fraction field
of \(O\). Integral closure of the DVR puts its cleared
value in \(O\). For each abelian label with an allowed
\(s\), use one such \(s\) for its fixed coefficient
functions and then take the maximum over the finitely
many labels of this row. This gives the exponent \(N_0\)
in the abelian part of the simultaneous specialization.
\label{cor:toric-twist-1}
\end{proof}

\paragraph{The local jet bound in period coordinates.}
Let \(\omega^{12}\) be the fixed lattice discriminant
and let \(g\) be the defining holomorphic lattice sum
before inserting the powers of the conjugate variables.
The field data here use \(E/\mathbb Q\) abelian,
\(E/F\) unramified at all finite places, and a CM type
containing the chosen image \(i\sqrt d\) of \(\sqrt{-d}\)
at each selected embedding. In particular
\([F:\mathbb Q]\ge2\). If \(F=\mathbb Q\), the
discriminant formula would give \(|D_E|=1\) for an
imaginary quadratic field. The condition on the real
degree applies at weight two as well.

\begin{proof}[Proof of Proposition~\ref{prop:toric-class-denominator}]
\proofstep{Integral idempotents and tangent directions.}
We give the integral comparisons on the lattice
\(E\cap y\widehat{\mathcal O}_{\mathfrak c}\) with
\(y=1\) above \(\ell\). Work over the valuation ring
\(R\) of a finite algebraic coefficient field in
the prescribed embeddings. We enlarge it finitely
to contain the CM embeddings and the good-reduction
model constructed below. The order action extends
to that abelian scheme by the extension property
for homomorphisms. Use \(v_\ell(\ell)=1\) in every
coefficient extension.
Use a generator of the fractional polarization
module locally at \(\ell\). If \(H\) denotes the integral de Rham lattice, perfection of the resulting alternating form means
\[
 H\xrightarrow{\;v\mapsto B(v,-)\;}H^\vee\ \text{is an isomorphism}.
\]

Indeed that local generator gives an isomorphism to
the dual at the coefficient prime, and the integral
de Rham duality comparison gives its perfect form.
The comparison with the ordinary polarization
pairing uses \(2B\), with \(B\) itself the perfect
integral pairing. The generic CM eigenspaces have
rank one, with Hodge subspace given by one
embedding from each conjugate pair. The perfect
pairing is between the opposite eigenspaces.

Let \(D\) be a scalar clearing every primitive CM
idempotent, with
\[
 v_\ell(D)=c_\ell.
\]
For a basis
\((x_j)_j\) of the maximal order and its trace-dual basis
\((x_j^\vee)_j\), the idempotent at an embedding \(\iota_i\)
is \(\sum_j\iota_i(x_j^\vee)x_j\) after coefficient
extension. Multiplication by an element of the different
clears the trace-dual coordinates. Galois uniformity gives
\[
 c_\ell=v_\ell(D_E)/[E:\mathbb Q]
\]
away from the conductor.
At a conductor prime the maximal order is unramified at the
odd prime \(\ell\), so those idempotents are integral there.
Multiplication by \(\ell^{a_0}\) puts their coefficients in
the conductor order. The extra order places are prime to
\(\ell\), so the same local order works in both comparisons.

Let \(e_i,f_i\) be the primitive generators of the two
opposite integral lines in the de Rham lattice and put
\(P_i=B(e_i,f_i)\). They are saturated line lattices.
Perfection of \(B\) supplies an integral \(x\) with
\(B(e_i,x)=1\). Write \(D E_{\bar i}x=s_i f_i\), where
\(E_{\bar i}\) is the idempotent on the opposite line and
\(s_i\in R\). Then
\[
      s_iP_i=B(e_i,D E_{\bar i}x)=D ,
      \qquad 0\le v_\ell(P_i)\le c_\ell .
\]
The equality uses the oppositely paired eigenspaces.
Multiplying the rank-one map \(e_i\mapsto f_i\) by \(D\)
gives an integral operator \(M_i\) killing other eigenlines,
as \(D E_i\) maps into the primitive lattice on the
line of \(e_i\), where \(E_i\) denotes that idempotent.
Each \(M_i\) commutes with the real order and is skew for \(B\).

If \(M=\sum_i t_i M_i\), then \(M^2=0\) and
\[
 B((1+M)u,(1+M)v)
  =B(u,v)+B(Mu,v)+B(u,Mv)+B(Mu,Mv)=B(u,v).
\]
The linear terms cancel by skewness and the last is
\(-B(u,M^2v)=0\). The same expansion applies to the other
polarization forms, since they are obtained by real
multiplication and the \(M_i\) commute with that multiplication.

\proofstep{Comparison with the reference CM periods.}
The primitive integral tangent eigenforms have
complex period scalars \(a_i\). For their comparison
with \(\omega\), let \(A_{\rm ref}\) be the elliptic curve
of the fixed quadratic lattice. Its CM \(j\)-invariant
is integral, so it has good reduction after finite
extension at \(\ell\). Choose an integral generator of
its tangent line. Its analytic coordinate \(t\) in the
given lattice has \(t^{12}\omega^{12}\) equal to the
discriminant of a good Weierstrass model. This is the
weight-twelve change of Weierstrass invariants, using
the full Weierstrass equation at two and three.
Thus \((t\omega)^{-1}\) is algebraic and integral at
the selected place.

Let
\[
 L_y=E\cap y\widehat{\mathcal O}_{\mathfrak c}.
\]

Since the local lattice at each place above \(\ell\)
is the conductor order, \(\ell^{a_0}\mathcal O_E\)
is contained in it locally. At the remaining places
one integer \(m_0\) prime to \(\ell\) clears the
necessary lattice index. Thus
\[
 n_0\mathcal O_E\subset L_y
\]
with
\(n_0=\ell^{a_0}m_0\). Choose \(\theta\in\mathcal O_E\)
separating the complex embeddings and take the
matrix with columns \(n_0(\varphi_i(\theta^j))_i\),
\(0\le j<f_0\).

For \(V=(n_0\varphi_i(\theta^j))_{i,j}\), the embeddings
give distinct values \(\varphi_i(\theta)\), so
\[
 \det V=n_0^{f_0}\prod_{i<i'}
               (\varphi_{i'}(\theta)-\varphi_i(\theta))\ne0 .
\]
The reference lattice consists of simultaneous embedded
quadratic integers. Thus column \(j\) maps such an
integer \(a\) to the lattice element \(n_0a\theta^j\)
of \(n_0\mathcal O_E\subseteq L_y\). The matrix defines an
isogeny from \(A_{\rm ref}^{f_0}\) to the CM variety.
After finite extension its finite generic kernel has
a finite flat schematic closure in the torsion group
of the good source. Quotient by this closure gives
the good target model and an integral isogeny.
The integral tangent vector on the factor in column
\(j=0\) has image with analytic coordinates
\(w_i^0=n_0t\) for all \(i\). Consequently
\[
 v_\ell((w_i^0\omega)^{-1})\ge-a_0.
\]

Evaluating the integral eigenform on this vector
gives \(a_iw_i^0\in R\). This implies the first
of the two bounds below. For the determinant,
the same cleared CM idempotents on the tangent
dual put \(D\) times the full lattice in the span
of the primitive eigenforms. Their wedge is
\(u_0\) times a determinant generator, where
\(0\le v_\ell(u_0)\le f_0c_\ell\) by Smith normal
form. The complex determinant scalar is
\[
 C=(\prod_i a_i)/u_0.
\]


For the lowering coordinates let \(w_i\) be the
complex tangent coordinate of \(f_i\) and choose
the integral unit \(r_i\) with \(P_i=r_i a_i w_i\).
This equality follows by using \(B(e_i,-)\) as
the primitive covector on the tangent quotient.
The Riemann form in the polarized uniformization
has real scale \(h_i\) on this coordinate.
If \(\psi\) is the complex polarization pairing
and \(\xi\) is the imaginary generator of the
polarization module locally at \(\ell\), the
comparison of alternating forms gives
\[
  2i\varphi_i(\xi)\psi(e_i,f_i)=h_i P_i .
\]
The local ideal of polarizations is dual to the
lattice under the trace pairing.

At the chosen
places, \(y=1\) and the extra level is a unit.
The local generator consequently satisfies
\(v_\ell(\varphi_i(\xi))\le0\).
Equivalently one can choose a local anti-invariant
inverse-different generator \(\xi'\) of nonpositive
valuation and use
\(m\xi'=a\xi\) with \(m\) prime to \(\ell\)
and \(a\in\mathcal O_F\) to compare the two generators.
The scale of the lowering direction and its auxiliary
period are
\[
  A_i=\pi h_i\frac{D}{\psi(e_i,f_i)}w_i^2,\qquad
  b_i=\frac{a_i A_i}{D}.
\]
In particular \(b_i/(a_i A_i)=D^{-1}\). Writing
\(\delta=\sqrt{-d}\) with \(\varphi_i(\delta)=i\sqrt d\),
the other period quotient is
\[
   (a_i b_i\sqrt d/(2\pi))^{-1}
       =\frac{r_i^2}{P_i\,\varphi_i(\xi\delta)} .
\]
All of these right-hand sides are algebraic.
Since \(v_\ell(\varphi_i(\delta))=\frac12\mathbf1_{\ell=d}\),
we obtain
\[
 v_\ell(a_i/\omega)\ge-a_0,\qquad
 v_\ell\bigl((a_i b_i\sqrt d/(2\pi))^{-1}\bigr)
     \ge-c_\ell-\tfrac12\mathbf1_{\ell=d}.
\]
The determinant and inverse-idempotent bounds just
proved supply the other two valuations.

\proofstep{The divided-power deformation.}
Use the divided-power thickening of the chosen abelian scheme. In horizontal coordinates the prescribed variation of its Hodge submodule is
\[
 \operatorname{Fil}^1_{\boldsymbol t}=(1+\sum_i t_iM_i)\operatorname{Fil}^1_{0},
\]
with the original filtration extended to the jet ring.
The real action and polarizations lift with the
filtration because the indicated operators commute
with that action and preserve the forms.
Here use the divided-power jet lifting theorem for the integral level.
Its base is a characteristic-zero local domain and its
thickening is the truncated divided-power jet ring,
with split augmentation and nilpotent divided-power
kernel. The input endomorphisms are admissible for the
real action, filtration and alternating forms, as are the
lowering operators above. The theorem gives a polarized
abelian scheme with real multiplication over that jet
ring, together with the pullback to the augmentation.
Under the horizontal de Rham identification, its
Hodge submodule is the span of the prescribed jet lifts
of the original filtration.

If \(s_{\det}\) is the determinant lift and \(s_0\) the given generator at the origin, its normalization in the horizontal identification is
\[
 s_{\det}\bmod(t_1,\ldots,t_d)=s_0,
\]
where \(d\) is the number of deformation parameters.
Complex base change identifies the integral jet and its lifted filtration with the analytic deformation used to calculate the periods.
On the reductions modulo powers of \(\ell\) the
underlying deformation comparison has \(\ell\) nilpotent,
as in Grothendieck--Messing theory.
\cite{messing}\label{cor:toric-twist-2}

Transport the periods horizontally in that family.
In the \(i\)-th real factor take complex combinations
of the parallel cycles with holomorphic coordinate one
and with zero holomorphic coordinate and antiholomorphic
coordinate one, respectively. The latter cycle
represents the Hodge direction at the origin.
After the lowering it lies in the annihilator of
the lifted holomorphic forms.

Let \(P_i^{\rm hol}\) be the nonzero period on the first cycle and \(\varepsilon\) the period pairing sign. Evaluation on the second cycle gives
\[
          P_i^{\rm hol}(\zeta_i-\varepsilon A_i t_i)=0 .
\]
Its constant term is nonzero, so \(P_i^{\rm hol}\) is a unit
in the complex jet ring. Hence \(\zeta_i=\varepsilon A_i t_i\) there.
The determinant scale on the horizontal frame is
\(C\). Indeed the parallel cycles with holomorphic
coordinates the standard basis are on the opposite
CM eigenlines, where the lowering operator is zero.
Evaluating the lifted determinant on their wedge
gives exactly the value at the origin.

\proofstep{Extracting the coefficient of the prescribed jet.}
Put \(g_\kappa=\partial^\kappa g(0)\) and
\[
 J_\kappa=g_\kappa\prod_i((b_i/a_i)^{\kappa_i}/a_i^{2s}).
\]

Let \(m=\ell^{c_E}\) and let \(R\) be the local ring
of algebraic coefficients for the integral CM model,
contained in the chosen valuation ring. It contains
the rational \(\ell\)-integers. The weight-\(2s\) Hilbert
Eisenstein form on each polarized component has
integral nonconstant Fourier coefficients in the
algebraic normalization. Its constant coefficient
is the corresponding partial zeta value divided
by \(2^{f_0}\). The partial zeta rationality theorem
applies since \(F\) is totally real and \(s\ge1\).
The integral model of the Hilbert moduli problem
and the \(q\)-expansion principle on all components
then gives the form multiplied by \(m\) over \(R\).
For weight two the Fourier statement uses the value
at zero of the continued family in real degree at
least two. \cite{deligne-pappas,serre-tate}\label{cor:toric-twist-3}

Use the divided-power algebra in the deformation parameters truncated at total degree
\[
 |\kappa|=\sum_i\kappa_i.
\]

The admitted de Rham operators of the integral CM
model determine its deformation over this algebra.
The lifted determinant generator reads the section
as a scalar in the divided-power basis, with
coefficient \(x\in R\) at the multi-index \(\kappa\).
After base change to \(\mathbb C\), the lowering
direction in the \(i\)-th coordinate has scale
\(\varepsilon A_i\), where \(\varepsilon\in\{1,-1\}\).
The determinant scale is \(C\). Comparing the divided
power coefficient after base change gives
\[
       C^{2s}x=m g_\kappa\prod_i(\varepsilon A_i)^{\kappa_i}.
\]

In particular there is no factorial in this coefficient
identity. The integral model comparison gives
\(v_\ell(C/\prod_i a_i)\ge-f_0c_\ell\) and
\(v_\ell(b_i/(a_i A_i))\ge-c_\ell\). Thus
\[
       J_\kappa=x m^{-1}
         \left(\frac{C}{\prod_i a_i}\right)^{2s}
         \prod_i\left(\frac{b_i}{\varepsilon a_i A_i}\right)^{\kappa_i}
\]
has valuation at least \(-c_E-2sf_0c_\ell-K_\kappa c_\ell\).
At zero depth this is \(-c_E\).
At positive depth the sum of the factorial valuations satisfies
\[
                  \sum_i v_\ell(\kappa_i!)\le K_\kappa/(\ell-1).
\]
Subtract this from the bound for \(v_\ell(J_\kappa)\) to obtain
\[
 v_\ell\!\left(J_\kappa/\prod_i\kappa_i!\right)
    \ge -c_E-2sf_0c_\ell
             -K_\kappa(c_\ell+1+\mathbf1_{\ell=2}) .
\]
\proofstep{Returning to the harmonic lattice sum.}
The normalized coordinate derivative for the harmonic
lattice sum is \((-1)^{K_\kappa}c_{2s}
\prod_i(2s)_{\kappa_i} S_y(0)\), where \(c_{2s}\) is
the defining lattice constant and \((2s)_a\) is the
rising factorial. This identity follows by differentiating
the continued defining sum. In the convergence region each
derivative increases the relevant denominator exponent by
one, supplying its successive factor, and continuity in
the exponent gives the formula at the evaluation point.
The definition of \(\Omega_*\) absorbs the same constant
and rising factorials. Multiplying \(J_\kappa\) by
\[
   \prod_i (a_i b_i\sqrt d/(2\pi))^{-\kappa_i}
                  (a_i/\omega)^{2s+2\kappa_i}
\]
leaves, up to the sign, \(S_y(0)/\Omega_*\).

In the positive-depth case use the divided-power jet
and also multiply by \(\prod_i\kappa_i!\). These
factorials are integral at every prime. The two period
factors cost \(K_\kappa(c_\ell+\tfrac12\mathbf1_{\ell=d})\)
and \(Wa_0\), respectively. At zero depth \(\ell\ne d\)
since the auxiliary field contains the imaginary quadratic
subfield of discriminant \(-d\). The same multiplication
then costs only \(Wa_0\). Adding these contributions
and taking the twelfth power proves the stated local
lattice estimate from the normalized jet inequalities.
The coordinate and jet comparisons use the fixed CM
type and base conductor. The auxiliary order and the
class label may vary with these data and the coefficient
prime fixed. \label{cor:toric-twist-4}
\end{proof}

\paragraph{Distribution at an added split place.}
Write \(q\) for the size of the residue field of a finite place \(v\) of the totally real field
that splits in the CM extension, and \(P,P^c\) for the places above it.
The local completed order at depth one consists of
pairs with equal residues. For a character unramified
at \(v\), let \(u\) be its value on a uniformizer at
\(P\), so its value at \(P^c\) is \(u^{-1}\). In the
absolutely convergent half-plane write \(X=q^{-z}\),
where \(z\) is the complex argument of the Hecke series.

\begin{lemma}[Distribution at a split finite place]\label{lem:toric-split-distribution}
With inverse Artin reciprocity on the ideal classes, removal of an
auxiliary pair multiplies a value unramified at the pair by
\[
       (1-uX)(1-u^{-1}X)+(q-1)X^2,\qquad X=q^{-z},
\]
where \(u\) is the corresponding twisted uniformizer value on
one factor. This relation holds on the remaining group-ring
coordinates with the same period normalization.
\end{lemma}

\begin{proof}
The two unit coordinates with equal residues give one
class modulo the order units. Pairs with both
coordinates divisible by a uniformizer give \(q-1\)
unit cosets for each pair of positive valuations.
There are no pairs with exactly one positive valuation.
The local sum is consequently
\[
     1+(q-1)\frac{X^2}{(1-uX)(1-u^{-1}X)} .
\]
Relative to the original two Euler factors this is the
polynomial
\[
          (1-uX)(1-u^{-1}X)+(q-1)X^2 .
\]
If the tested character is ramified at this place,
only the simultaneous unit class contributes. The other
terms sum to zero because the character on the ratio of
residue units is nontrivial. The usual local factor of
the Hecke series there is one as well. The exponent-one
order is essential to the unit coset count.

\proofstep{The unit group of the order.}
To justify applying this computation to the lattice sum,
we identify the global units of the order precisely.
The base conductor is nonzero and proper and has an odd
prime divisor of residual characteristic \(p\) unramified
in \(E\) over \(\mathbb Q\). Let \(\mathfrak P\) be a prime of
\(E\) above that divisor. If \(u\in\mathcal O_{\mathfrak c}^\times\),
then \(u,\bar u\) are algebraic units and
\[
 u-\bar u\in\mathfrak c\mathcal O_E\subseteq\mathfrak P,
\]
since the image of the order modulo the conductor is real.
All complex conjugates of \(u/\bar u\) have absolute value one.
An algebraic integer with that property is a root of unity,
as the monic polynomials of all its powers have bounded integral
coefficients. At \(\mathfrak P\) the root lies in the principal unit
group. Here the normalized local valuation satisfies \(v(p)=1\).

If \(v(t)\ge1\) and \(t\ne0\), then \(v(pt)=1+v(t)\). The
intermediate terms \(\binom pj t^j\), \(2\le j<p\), have valuation
at least \(1+2v(t)\), and the final term has valuation \(pv(t)>1+v(t)\)
because \(p\) is odd. Thus
\[
 v((1+t)^p-1)=v(pt)=1+v(t).
\] For an exponent \(a\) prime to \(p\) the linear term
gives \(v((1+t)^a-1)=v(t)\). For any positive exponent
\(m=p^j a\) with \((a,p)=1\), iteration therefore gives
\[
                  v((1+t)^m-1)=v(t)+j<\infty .
\]
This is incompatible with \(m\)-torsion.
Thus the principal unit group has no nonidentity torsion, and
\(u=\bar u\). Both \(u\) and its inverse belong to \(\mathcal O_F\),
so
\[
 \mathcal O_{\mathfrak c}^\times=\mathcal O_F^\times.
\]


With this equality,
in the convergence half-plane group the nonzero lattice
vectors by the corresponding integral finite ideles
modulo the completed order units. Dividing each
idele into a class representative and an element of
\(E^\times\) identifies the fiber with the lattice vectors
modulo the global order units. Homogeneity identifies
its summand with the Hecke weight. The local
order conditions and the multiplicativity of the weight
give the Euler product. At a place of the base conductor
the character has the prescribed exact conductor.
In its nonunit layers the sum over the last available
residue unit quotient vanishes by that conductor, leaving
the primitive local factor. At all other unchanged
places the full order gives the usual geometric series.
The equality of the continued sums follows from the
identity theorem, and at the desired value one
substitutes
\[
 X=q^{-s}.
\]


\proofstep{The sign in the inverse-class map.}
Pair removal will compare coefficients at a split prime through their ray classes. The transfer must therefore use the same convention for inversion as the distribution relation.

For removal of one pair of auxiliary places, evaluate
the abelian group-ring data on any character of the
remaining Artin labels. Such a character is unramified
at the removed place. Summing the removed coordinate
amounts to forgetting it in the inverse-class map.
A uniformizer at either place above \(v\) acts on the
remaining labels through its inverse Artin class.
Define \(Y\) to be the inverse of the image of \(\chi(P)\)
in the coefficient field.

With this reciprocity convention the unramified twisting factor at \(P\) is
the value of the remaining character on its Frobenius.
Thus evaluation at a remaining character \(\psi\) gives
\[
 u_\psi=Y^{-1}\psi(\operatorname{Frob}_P),\qquad
 (1-u_\psi X)(1-u_\psi^{-1}X)+(q-1)X^2 .
\]
The second local root is \(u_\psi^{-1}\).
The polynomial gives the same
multiplicative relation at every tested character, so
finite Fourier inversion proves the group-ring
distribution. With all auxiliary places removed, the
sum of the base coefficients is a nonzero scalar times
the untwisted Hecke \(L(s,\chi)\). The character is unitary
and is ramified at the base conductor, so the classical
nonvanishing theorem for unitary Hecke \(L\)-functions
on \(\RePart z\ge1\), away from the pole of a norm
character, applies. Both the coefficient bound and this
distribution refer to the same inverse-class map and
the same period. \label{cor:toric-twist-5}
\end{proof}
